% Created 2018-02-02 Fri 12:29
% Intended LaTeX compiler: pdflatex
\documentclass[11pt]{article}
\usepackage[utf8]{inputenc}
\usepackage[T1]{fontenc}
\usepackage{graphicx}
\usepackage{grffile}
\usepackage{longtable}
\usepackage{wrapfig}
\usepackage{rotating}
\usepackage[normalem]{ulem}
\usepackage{amsmath}
\usepackage{textcomp}
\usepackage{amssymb}
\usepackage{capt-of}
\usepackage{hyperref}
\author{Monib Ahmed}
\date{\today}
\title{Simple Arithmetic Unit}
\hypersetup{
 pdfauthor={Monib Ahmed},
 pdftitle={Simple Arithmetic Unit},
 pdfkeywords={},
 pdfsubject={},
 pdfcreator={Emacs 25.3.1 (Org mode 9.1.5)}, 
 pdflang={English}}
\begin{document}

\maketitle
\tableofcontents


\section{Introduction}
\label{sec:org1f9f06a}
This is a simple AU (SAU) unit that is specifically used to walked
through the Python to Hardware Design flow. 

\section{Operations}
\label{sec:org5b3a8f7}
\begin{center}
\begin{tabular}{rl}
Opcode & Operation\\
2-bits & \\
\hline
000 & NOP\\
001 & ADD\\
010 & SUB\\
011 & SHIFTL\\
100 & SHIFTR\\
101 & AND\\
110 & OR\\
111 & NOP\\
\end{tabular}
\end{center}

\section{Interface}
\label{sec:org892b1e7}
\begin{center}
\begin{tabular}{lll}
Port Name & Direction & Description\\
\hline
clk & input & Clock\\
rstb & input & Reset\\
opcode[2:0] & input & Opcode\\
A[3:0] & input & Operand A\\
B[3:0] & input & Operand B\\
C[4:0] & output & Result\\
\end{tabular}
\end{center}
\end{document}
